% ============================================================
% Asilomar 1-Page Extended Abstract — Polished Draft
% Compile with: pdflatex asilomar_support_bottleneck_polished.tex
% ============================================================

\documentclass[conference,10pt]{IEEEtran}

\usepackage{amsmath, amssymb}
\usepackage{booktabs}
\usepackage{cite}

\title{Support Before Amplitude: A Bottleneck in Cross-Operator Sparse Recovery}

\author{
\IEEEauthorblockN{Di An}
\IEEEauthorblockA{Johns Hopkins University\\
Email: dan5@jhu.edu}
\and
\IEEEauthorblockN{Dylan Poppert}
\IEEEauthorblockA{Johns Hopkins University\\
Email: dpopper2@jhu.edu}
\and
\IEEEauthorblockN{Trac D. Tran}
\IEEEauthorblockA{Johns Hopkins University\\
Email: trac@jhu.edu}
}

\begin{document}
\maketitle

% \begin{abstract}
% We study zero-shot transfer in sparse recovery when the sensing operator
% changes between training and deployment. In the setting $n=256$,
% $m=128$, and $k=25$, we train on Fourier measurements and evaluate on
% Gaussian measurements. A simple parameter-free pipeline, consisting of
% ISTA, naive top-$k$ support selection, and support-restricted least
% squares, achieves Gaussian-test NRMSE $0.257$, improving over zero-shot
% LISTA trained on the same Fourier data ($0.327$) by $21\%$ relative
% NRMSE. Surprisingly, two learned coordinate-wise support detectors of
% different capacity, both trained on ISTA-derived iterate features, match
% but do not improve over naive magnitude ranking, with NRMSE differences
% within $0.001$ across several $(T,\lambda)$ settings.

% These results suggest that the cross-operator transfer failure is not
% primarily an amplitude-estimation problem: with oracle support,
% support-restricted least squares gives nearly exact recovery. Instead,
% the bottleneck is support identification under operator shift. We
% attribute the plateau of learned detectors to their iterate-only feature
% interface, which exposes source-operator magnitude statistics but little
% explicit information about the geometry of the deployment operator. This
% diagnosis motivates an operator-aware support detector that augments each
% coordinate with residual correlations, column norms, and Gram/coherence
% features derived from the test-time sensing matrix. The resulting
% decomposition preserves the classical support-then-amplitude structure
% while asking whether lightweight operator-aware features can close the
% gap to target-operator fine-tuning without requiring target labels.
% \end{abstract}

\begin{abstract}
We study zero-shot sparse recovery under operator shift. On
Fourier $\to$ Gaussian transfer ($n{=}256$, $m{=}128$, $k{=}25$),
a parameter-free pipeline of ISTA, naive top-$k$ support selection,
and support-restricted least squares achieves NRMSE $0.257$, improving
over zero-shot LISTA ($0.327$) by $21\%$ relative. Two learned
coordinate-wise detectors of differing capacity, both operating on
iterate-derived features, tie naive magnitude ranking within $0.001$
NRMSE. Because oracle support yields near-exact recovery, the binding
constraint is support identification, not amplitude estimation. We
trace the detector plateau to an iterate-only feature interface and
propose operator-aware coordinate features---residual correlations and
operator Gram statistics computed from the deployment matrix---as a
minimal learned module for transfer without target labels.
\end{abstract}

\begin{IEEEkeywords}
Sparse recovery, compressed sensing, learned iterative reconstruction,
LISTA, transfer learning, support detection.
\end{IEEEkeywords}

% ------------------------------------------------------------
\section{Problem Setup}
% ------------------------------------------------------------

We consider sparse recovery from linear measurements
\begin{equation}
    y = A x + \varepsilon,
\end{equation}
where $x \in \mathbb{R}^{n}$ is $k$-sparse and
$A \in \mathbb{R}^{m \times n}$ is the sensing matrix. Given a fixed
operator $A$, we first compute an ISTA iterate
\begin{equation}
    x^{(T)} = \mathrm{ISTA}_{T}(y; A, \lambda).
\end{equation}
We then compare four recovery strategies under operator shift. All
learned models are trained on $4{,}000$ Fourier-measurement signals and
evaluated zero-shot on Gaussian measurements: (i) raw ISTA; (ii) naive
top-$k$ support selection on $|x^{(T)}|$ followed by support-restricted
least squares; (iii) a learned coordinate-wise support detector operating
on iterate-derived features, again followed by least squares; and
(iv) end-to-end LISTA.

This setting separates two questions that are often entangled in learned
sparse recovery: whether the method identifies the correct support, and
whether it estimates the amplitudes accurately once the support is known.

% ------------------------------------------------------------
\section{Main Result}
% ------------------------------------------------------------

\begin{table}[t]
\centering
\caption{Zero-shot Fourier $\to$ Gaussian transfer at $T=30$,
$\lambda=0.05$, $n=256$, $m=128$, and $k=25$.}
\label{tab:main}
\vspace{-4pt}
\begin{tabular}{lc}
\toprule
\textbf{Method} & \textbf{Gaussian NRMSE} \\
\midrule
Raw ISTA                                  & 0.4465 \\
Naive top-$k$ + LS                        & \textbf{0.2574} \\
Learned detector + LS                     & 0.2578 \\
LISTA, zero-shot                          & 0.3267 \\
LISTA, $50$ Gaussian fine-tuning samples  & 0.2551 \\
Oracle support + LS                       & $\approx 0$ \\
\bottomrule
\end{tabular}
\vspace{-6pt}
\end{table}

Table~\ref{tab:main} shows that the parameter-free support-then-amplitude
pipeline outperforms zero-shot LISTA by $21\%$ relative NRMSE. At the
same time, learned support detectors operating on ISTA-derived features
track naive top-$k$ selection almost exactly: across
$(T,\lambda) \in \{(10,0.05),(30,0.05),(30,0.01)\}$ and across two
detector architectures, their NRMSE differs from naive ranking by less
than $0.001$. With approximately $50$ Gaussian fine-tuning samples, LISTA
closes the gap, indicating that the zero-shot advantage is most relevant
in the small-target-data regime.

% ------------------------------------------------------------
\section{Diagnosis: Support, Not Amplitude}
% ------------------------------------------------------------

The oracle-support result isolates the source of error. Once the correct
support is known, support-restricted least squares gives nearly exact
recovery. Therefore, the gap between the practical methods and oracle
recovery is mainly caused by support misidentification rather than by
inaccurate coefficient estimation.

The plateau of learned detectors suggests that the feature interface is
too restrictive. The detector input at coordinate $j$ is derived from
$x^{(T)}_j$ and related magnitude statistics. However, $x^{(T)}$ is itself
a nonlinear product of the ISTA dynamics and the sensing operator. Under a
Fourier $\to$ Gaussian shift, these iterate-only features preserve the
same coordinate-ranking information already exploited by naive top-$k$
selection, but they do not explicitly encode the geometry of the
deployment operator. This explains why the learned detector can match the
magnitude baseline yet fails to improve over it.

% ------------------------------------------------------------
\section{Toward Operator-Aware Support Detection}
% ------------------------------------------------------------

The diagnosis suggests that support detectors should not be restricted to
iterate-derived magnitude features. For each coordinate $j$, we propose
augmenting the detector input with operator-aware features computed at
test time:
\begin{equation}
\phi_j =
\left[
|x^{(T)}_j|,
\; A_{:,j}^{\top}(y-Ax^{(T)}),
\; \|A_{:,j}\|_2,
\; (A^{\top}A)_{jj},
\; \max_{i\neq j}|A_{:,i}^{\top}A_{:,j}|
\right].
\end{equation}
The residual correlation term measures whether coordinate $j$ can still
explain the current measurement residual, while the column norm, Gram
diagonal, and coherence terms describe the local identifiability of that
coordinate under the deployment operator. A lightweight coordinate-wise
classifier can then estimate support probabilities from $\phi_j$, after
which amplitudes are computed by least squares on the selected support.

This design reintroduces the operator-dependent information exploited by
unrolled methods such as LISTA, but confines learning to the support
decision stage. The key empirical question is whether this operator-aware
interface can improve zero-shot cross-operator support recovery and
approach the performance of target-operator fine-tuning without using
target labels.

% ------------------------------------------------------------
\section{Conclusion}
% ------------------------------------------------------------

Our results identify a support bottleneck in zero-shot cross-operator
sparse recovery. In the studied Fourier $\to$ Gaussian setting, a simple
ISTA + top-$k$ + least-squares pipeline outperforms zero-shot LISTA, while
learned iterate-only support detectors fail to exceed naive magnitude
ranking. These findings motivate operator-aware support detection as a
minimal learned module for improving transfer: rather than learning the
entire inverse map, the model learns only the support decision using both
iterate state and test-time operator geometry.

\end{document}
